\documentclass[../structural_economics_book.tex]{subfiles}
\begin{document}
\ifSubfilesClassLoaded{\setcounter{chapter}{26}}{}
\chapter{Algorithmic $W^I$}\label{ch:27}
\chaptersubtitle{When the Machine Rewires Who Sees What}
\begin{quote}
\textit{``The algorithms that determine what you see are not neutral. They are structural interventions on the information network --- operators on $W^I$ with objectives that may have nothing to do with your welfare or your society's development.''}\\
\hfill --- This book
\end{quote}

\vspace{1em}

Every chapter until now has treated $W^I$ as shaped by human institutions: media organizations, schools, social networks, community structures. This chapter confronts the twenty-first century's structural novelty: $W^I$ is now substantially designed by algorithms.

Recommendation systems on Facebook, TikTok, YouTube, Google, and WhatsApp determine what billions of people see, whom they hear, and which signals enter their formation equation. These algorithms are operators on $W^I$ --- they add, remove, and reweight information edges in real time. They are the most powerful $W^I$ operators ever deployed, operating at a scale and speed that no human institution has ever achieved.

The question for structural economics: do these algorithms increase or decrease $N_{\text{eff}}$? Do they break assortative mixing or reinforce it? Do they connect the quinoa farmer in Potos\'{i} to the engineer in Santa Cruz, or do they lock each into a personalized information silo?

The answer: it depends entirely on the algorithm's objective function. And that objective function is not set by the farmer, the engineer, or the Bolivian state. It is set by firms in Menlo Park, Beijing, and Mountain View.


%% =====================================================================
\section{The Algorithm as an Operator on $W^I$}
\label{ch27:sec:operator}
%% =====================================================================

\subsection{Formalization}

\begin{definition}[Algorithmic $W^I$ Operator]\label{ch27:def:algorithmic}
A \textbf{recommendation algorithm} is a function $\mathcal{A}: (W^I, \mathbf{h}) \to W^{I'}$ that takes the current information network $W^I$ and a vector of engagement histories $\mathbf{h} = (h_1, \ldots, h_N)$ and outputs a modified information network $W^{I'}$:
\begin{equation}\label{ch27:eq:algorithmic}
W^{I'}_{ij} = \mathcal{A}_{ij}(W^I, h_i, h_j, \boldsymbol{\eta}),
\end{equation}
where $\boldsymbol{\eta}$ are the algorithm's parameters (learned from data to maximize an objective).

The algorithm modifies $W^I$ by choosing which content from source $j$ to show user $i$. If the algorithm shows $i$ more content from $j$: $W^{I'}_{ij} > W^I_{ij}$ (the edge strengthens). If it suppresses content from $j$: $W^{I'}_{ij} < W^I_{ij}$ (the edge weakens).
\end{definition}

\begin{proposition}[The Algorithm's Objective Determines $W^I$'s Structure]\label{ch27:prop:objective}
Three objective functions produce qualitatively different $W^I$ structures:

\textbf{Objective 1: Maximize engagement} (the standard commercial objective). The algorithm shows $i$ content that maximizes $i$'s time on platform. Engagement is maximized by content that is emotionally arousing, confirms existing beliefs, and comes from sources similar to $i$. This produces \textbf{assortative amplification}: $W^{I'}_{ij}$ is large when $i$ and $j$ are similar (same demographic, same beliefs, same consumption pattern). The algorithm reinforces the existing block structure of $W^I$ and may decrease $N_{\text{eff}}$.

\textbf{Objective 2: Maximize information accuracy} (a hypothetical public-interest objective). The algorithm shows $i$ content from the most accurate and diverse sources. This produces \textbf{bridge-building}: $W^{I'}_{ij}$ is large when $j$ has information that $i$ lacks (cross-block edges). The algorithm breaks assortative mixing and increases $N_{\text{eff}}$.

\textbf{Objective 3: Maximize advertiser revenue} (a variant of engagement). Similar to Objective 1 but weighted by the commercial value of the user's attention. This produces \textbf{stratified exposure}: high-value users (urban, educated, high-income) see diverse, high-quality content; low-value users (rural, less-educated, low-income) see lower-quality, more assortative content. The algorithm increases $N_{\text{eff}}$ for the elite and decreases it for the periphery.
\end{proposition}

\begin{intuition}[The algorithm IS the new organic intellectual]
In Gramsci's framework (Chapter~24), the ``organic intellectual'' is the agent with high $\pi^I$ who produces the narratives that sustain the hegemony. In the algorithmic era, the organic intellectual is the recommendation system: it determines which narratives reach which audiences, with what frequency, and in what context. The algorithm has $\pi^I$ that dwarfs any human journalist, professor, or politician --- it reaches billions of users simultaneously and personalizes the message for each.

The difference: Gramsci's organic intellectual had an ideology (a preferred $W^I$ structure). The algorithm has an objective function (maximize engagement, revenue, or some other metric). The objective function may or may not align with any political project, any development goal, or any welfare criterion from Chapter~12. It is structural power without structural intent.
\end{intuition}


%% =====================================================================
\section{Three Effects on Structural Economics}
\label{ch27:sec:effects}
%% =====================================================================

\subsection{Effect 1: Echo Chambers and the Formation Trap}

\begin{theorem}[Algorithmic Amplification of Assortative Mixing]\label{ch27:thm:echo}
If the algorithm maximizes engagement and engagement is increasing in belief-confirmatory content, then:
\begin{equation}\label{ch27:eq:echo}
\lambda_k^{I'}(\mathcal{A}) \leq \lambda_k^I \quad \text{for all attributes } k,
\end{equation}
where $\lambda_k^{I'}$ is the cross-group mixing parameter after algorithmic filtering. The algorithm reduces cross-group information exposure, making $W^I$ more block-diagonal.

Consequence for the formation model: the self-confirming trap (Theorem~10.3) becomes deeper. The aspiration gap between groups widens because each group sees a more biased sample of returns. The $\Psi$ ratio increases (the formation channel becomes more dominant). The welfare cost of discrimination (Chapter~18) rises.
\end{theorem}

\begin{example}[Bolivia: Facebook and the Urban-Rural Information Divide]\label{ch27:ex:facebook}
Bolivia has 7 million Facebook users (2023). The recommendation algorithm shows users content similar to what they have previously engaged with. Urban Bolivians engage with commercial media, political commentary, and global content. Rural Bolivians engage with local community posts, agricultural content, and regional news.

The algorithm amplifies this difference: urban users see progressively more urban content; rural users see progressively more rural content. The cross-group mixing parameter $\lambda_{\text{urban-rural}}^I$, already low ($\approx 0.22$, Chapter~21), may decrease further under algorithmic filtering.

Estimated effect: $\Delta\lambda_{\text{urban-rural}}^I \approx -0.05$ to $-0.10$ (the algorithm reduces cross-group exposure by 20--45\% relative to a chronological feed). The formation model predicts: the aspiration gap between urban and rural populations widens by $\mu/(1-\rho) \times \Delta\lambda \times \Delta\bar{r} \approx 0.03$--$0.06$ preference units per generation.
\end{example}


\subsection{Effect 2: Epistemic Power Concentration}

\begin{proposition}[Algorithmic Concentration of $\boldsymbol{\pi}^I$]\label{ch27:prop:concentration}
Engagement-maximizing algorithms produce \textbf{winner-take-all dynamics} in $W^I$: a small number of viral content producers capture a large share of attention, while the vast majority of voices are suppressed. The Perron eigenvector $\boldsymbol{\pi}^I$ becomes more concentrated:
\begin{equation}\label{ch27:eq:concentration}
N_{\text{eff}}(W^{I'}) \leq N_{\text{eff}}(W^I).
\end{equation}

The mechanism: the algorithm amplifies content with high initial engagement (likes, shares, comments). Content from well-connected, emotionally compelling sources gets shown to more users, generating more engagement, generating more amplification. Content from less-connected sources (community radio, local journalists, indigenous-language media) gets suppressed because it generates less engagement.

The consequence: Bolivia's $\pi^I_{\text{community}}$ (already 0.05, Chapter~21) may decline further as algorithmic curation redirects attention to high-engagement commercial content. The $N_{\text{eff}}^I$ drops from 3.4 toward 2.5--3.0.
\end{proposition}


\subsection{Effect 3: Cross-Border Information Integration}

\begin{proposition}[The Paradox of Algorithmic Globalization]\label{ch27:prop:paradox}
Algorithms also create cross-border information edges that did not previously exist. A TikTok user in El Alto can watch a Korean entrepreneur, a Brazilian farmer, or an Indian programmer. The algorithm sometimes shows content from radically different contexts, creating $W^I$ edges that span the global product space.

This is Sunkel's integration/disintegration (Chapter~13) applied to $W^I$: the algorithm simultaneously integrates globally (showing diverse international content to some users) and fragments domestically (personalizing feeds within national populations). The net effect on domestic $N_{\text{eff}}$ is ambiguous:
\begin{equation}\label{ch27:eq:paradox}
\Delta N_{\text{eff}}^I = \underbrace{\Delta N_{\text{eff}}^{\text{global}}}_{\text{positive (new sources)}} + \underbrace{\Delta N_{\text{eff}}^{\text{domestic}}}_{\text{negative (echo chambers)}}.
\end{equation}

For urban educated Bolivians: the global effect may dominate (exposure to international entrepreneurship, technology, ideas). For rural indigenous Bolivians: the domestic effect may dominate (locked into algorithmically reinforced local content). The algorithm increases $N_{\text{eff}}$ for the elite and decreases it for the periphery --- amplifying the compound disadvantage of Chapter~19.
\end{proposition}


%% =====================================================================
\section{The Regulatory Response: $W^I$ as a Public Good}
\label{ch27:sec:regulation}
%% =====================================================================

\begin{definition}[Information Network as a Public Good]\label{ch27:def:public_good}
The structure of $W^I$ --- who sees what --- is a \textbf{public good} in the structural model: it determines the formation of preferences for the entire population, which determines the welfare loss (Chapter~2), the formation multiplier (Chapter~10), and the political equilibrium (Chapter~24). No individual user can choose the aggregate $W^I$ structure; it is the emergent outcome of billions of algorithmic decisions.
\end{definition}

\begin{proposition}[Four Regulatory Instruments]\label{ch27:prop:regulation}
\begin{enumerate}
\item \textbf{Transparency requirements}: Mandate disclosure of the algorithmic objective function and its effect on $W^I$ structure ($\Delta\lambda_k$, $\Delta N_{\text{eff}}$). If users and regulators can observe what the algorithm does to $W^I$, they can demand changes.

\item \textbf{Diversity mandates}: Require recommendation algorithms to maintain a minimum cross-group exposure ratio: $\lambda_k^{I'} \geq \lambda_k^{\min}$ for all attribute groups $k$. This prevents the algorithm from amplifying assortative mixing below a threshold.

\item \textbf{Public-interest algorithms}: Develop and deploy alternative recommendation systems with Objective~2 (maximize information accuracy/diversity) rather than Objective~1 (maximize engagement). These could be operated by public broadcasters, educational institutions, or platform cooperatives.

\item \textbf{$N_{\text{eff}}$ monitoring}: Track $N_{\text{eff}}(W^I)$ over time (Chapter~21) and treat declines as a regulatory trigger. If the algorithmically-mediated $N_{\text{eff}}$ falls below a threshold, regulatory intervention is warranted.
\end{enumerate}
\end{proposition}

\begin{keyinsight}
The algorithmic $W^I$ is the most important structural development of the twenty-first century. For the first time in history, the information network that shapes preference formation is designed by a small number of firms, operating at global scale, with objectives that are not aligned with any development goal.

The structural economics framework provides the language: the algorithm is an operator on $W^I$. Its effect is measurable ($\Delta\lambda_k$, $\Delta N_{\text{eff}}$, $\Delta\Psi$). Its welfare cost is computable (Chapter~18's welfare loss formula, amplified by $1 + \Psi$). And its regulation is an operator design problem (Chapter~15): choose the algorithmic parameters $\boldsymbol{\eta}$ that maximize $N_{\text{eff}}$ subject to commercial viability.

The framework transforms the debate from ``social media is good/bad'' to: ``what is the algorithm's effect on $\lambda_k$, $N_{\text{eff}}$, and $\Psi$, and does that effect move the structural trajectory toward or away from the growth-with-equity zone of Chapter~26?''
\end{keyinsight}


%% =====================================================================
\section{AI and the Future of Structural Analysis}
\label{ch27:sec:future}
%% =====================================================================

\subsection{AI as Measurement Tool}

\begin{proposition}[AI for Structural Diagnostics]\label{ch27:prop:ai_measurement}
Machine learning techniques can improve the measurement of all three weight matrices:
\begin{enumerate}
\item $W^P$: Satellite imagery + NLP on firm registries can estimate production linkages in real time, filling gaps between benchmark I-O tables (Chapter~20).
\item $W^I$: Natural language processing of social media content can estimate media overlap ($W^I_{\text{media}}$, Chapter~21) at daily frequency, tracking $N_{\text{eff}}$ in real time.
\item $W^U$: Sentiment analysis of social media posts can estimate solidarity patterns ($W^U_{\text{digital}}$, Chapter~22), supplementing survey-based measures.
\item The structural trajectory (Chapter~23) can be computed in near-real-time rather than at annual frequency.
\end{enumerate}
\end{proposition}


\subsection{AI as Structural Operator}

\begin{proposition}[AI-Powered Structural Intervention]\label{ch27:prop:ai_operator}
AI can also serve as a deliberate operator on $W^I$ for development purposes:
\begin{enumerate}
\item \textbf{Personalized mentoring at scale}: AI systems that connect aspiring entrepreneurs in the traditional sector with role models, market information, and technical knowledge from the modern sector. This is the $\mathcal{L}_{\text{info}}$ operator (Chapter~15) deployed through an AI system rather than a human program.
\item \textbf{Adaptive education}: AI tutoring systems that adjust content to the learner's context while exposing them to cross-group knowledge ($\Delta\lambda^I > 0$).
\item \textbf{Market information}: AI systems that provide real-time price and demand information to rural producers, creating $W^I$ edges from the traditional sector to the market that the assortative structure suppresses.
\end{enumerate}
\end{proposition}

\begin{prebisch}[The new center-periphery]
The algorithmic $W^I$ creates a new center-periphery structure that operates alongside the trade-based $W^P$:

\textbf{The algorithmic center}: Firms in Silicon Valley, Shenzhen, and a few other nodes control the parameters $\boldsymbol{\eta}$ of the global $W^I$. Their $\pi^I$ is enormous: a single change to Facebook's algorithm can shift the formation signals received by 3 billion people.

\textbf{The algorithmic periphery}: Every other country is a consumer of algorithmic $W^I$ designed elsewhere. Bolivia's $W^I$ is shaped by algorithms trained on US engagement data, optimized for US advertising markets, and deployed globally without adaptation to Bolivian development needs.

This is Prebisch's center-periphery thesis applied to $W^I$: the center produces the information infrastructure; the periphery consumes it. The ``terms of trade'' in $W^I$ are the exchange of attention (from the periphery) for algorithmic curation (from the center). The formation signals that Bolivian youth receive are filtered through an objective function designed in Menlo Park.

The structuralist response: just as Prebisch advocated for ISI to reduce $W^P$ dependency, the algorithmic era requires ``information substitution'' --- the development of domestic or regional algorithmic capacities that can design $W^I$ operators aligned with national development objectives. This is the information-age equivalent of building the capital goods sector: building the algorithmic infrastructure that shapes your own $W^I$.
\end{prebisch}


%% =====================================================================
\section{Summary}
\label{ch27:sec:summary}
%% =====================================================================

\begin{enumerate}
\item Recommendation algorithms are \textbf{operators on $W^I$}: they determine who sees what, with effects measurable as $\Delta\lambda_k$, $\Delta N_{\text{eff}}$, and $\Delta\Psi$.

\item \textbf{Engagement-maximizing algorithms} amplify assortative mixing ($\lambda_k \downarrow$), concentrate epistemic power ($N_{\text{eff}} \downarrow$), and deepen the formation trap ($\Psi \uparrow$). The welfare cost is computable.

\item The \textbf{paradox}: algorithms simultaneously integrate globally (cross-border content) and fragment domestically (echo chambers). The net effect differs by social position: elite $N_{\text{eff}}$ rises; peripheral $N_{\text{eff}}$ falls.

\item \textbf{Regulation}: transparency, diversity mandates, public-interest algorithms, $N_{\text{eff}}$ monitoring. The framework transforms the debate from ``good/bad'' to measurable structural effects.

\item \textbf{AI as tool}: machine learning can improve measurement of all three $W$'s in near-real-time. AI can also serve as a development operator (personalized mentoring, adaptive education, market information).

\item The \textbf{algorithmic center-periphery}: Silicon Valley designs the global $W^I$; the periphery consumes it. Information sovereignty requires building domestic algorithmic capacity --- the information-age capital goods sector.
\end{enumerate}


%% =====================================================================
\section{Exercises}
\label{ch27:sec:exercises}
%% =====================================================================

\begin{enumerate}[label=\textbf{27.\arabic*}, leftmargin=3em]

\item \textbf{Algorithmic $\Delta\lambda$.} For a $4 \times 4$ $W^I$ (urban men, urban women, rural men, rural women): simulate the effect of an engagement-maximizing algorithm that increases within-group exposure by 20\% and decreases cross-group exposure by 15\%. Compute $\Delta\lambda_k$ for gender and geography. Compute $\Delta N_{\text{eff}}$.

\item \textbf{The echo chamber theorem.} Prove Theorem~\ref{ch27:thm:echo}: if engagement is maximized by belief-confirmatory content, show that $\lambda_k^{I'} \leq \lambda_k^I$ under the standard collaborative-filtering algorithm.

\item \textbf{Bolivia's digital divide.} Using internet penetration data by department and demographic group: estimate the fraction of Bolivia's population whose $W^I$ is algorithmically mediated (vs.\ traditional media/word-of-mouth). How does this fraction correlate with $N_{\text{eff}}$ estimates from Chapter~21?

\item \textbf{Objective function design.} Propose an alternative algorithmic objective function for Bolivia that maximizes $N_{\text{eff}}$ subject to a minimum engagement constraint. Write the optimization problem formally. What trade-offs does it involve?

\item \textbf{AI mentoring simulation.} Design an AI-powered mentoring system for rural Bolivian entrepreneurs. Specify: (a) the $\Delta W^I$ edges it creates, (b) the expected $\Delta\lambda^I$ and $\Delta N_{\text{eff}}$, (c) the formation model prediction for $\Delta\theta^*$ after one generation.

\item \textbf{Algorithmic center-periphery.} Using Alexa/SimilarWeb data: estimate the fraction of Bolivia's internet traffic that goes to US-based platforms (Google, Facebook, YouTube, TikTok) vs.\ domestic content. What fraction of Bolivia's $W^I$ is algorithmically designed abroad?

\item \textbf{The policy frontier.} Compare four regulatory regimes: (a) no regulation (current), (b) transparency only, (c) diversity mandates ($\lambda_k^{\min} = 0.3$), (d) public-interest algorithm. For each: estimate the effect on $N_{\text{eff}}$, the political feasibility, and the cost. Which regime best advances Bolivia's structural trajectory?
\end{enumerate}


%% =====================================================================
\newpage
\section*{Mathematical Appendices}
\addcontentsline{toc}{section}{Mathematical Appendices}
%% =====================================================================


\subsection*{Appendix 27A: Recommendation Algorithms as Matrix Operators}
\addcontentsline{toc}{subsection}{Appendix 27A: Algorithms as Operators}
\label{ch27:app:algorithms}

\paragraph{Collaborative filtering.} The standard recommendation algorithm estimates a user-content affinity matrix $R \approx UV'$, where $U$ is a user-factor matrix and $V$ is a content-factor matrix (matrix factorization, analogous to the SVD of the specialization matrix in Chapter~16). Content shown to user $i$ is drawn from the rows of $V$ with highest inner product $\mathbf{u}_i' \mathbf{v}_j$.

\paragraph{Effect on $W^I$.} The user-factor matrix $U$ clusters users by ``type'' (implicit demographic, behavioral, or preference clusters). Content from source $j$ is shown to user $i$ only if $\mathbf{u}_i \approx \mathbf{u}_j$ (similar type). This produces assortative $W^I$: users see content from similar users. The mixing parameter is $\lambda_k^{I'} \approx \exp(-d(\bar{\mathbf{u}}_g, \bar{\mathbf{u}}_{g'}))$, where $d$ is the distance between group-average user factors. When groups are distinct in the factor space (different demographics, different languages, different consumption patterns): $d$ is large and $\lambda^{I'}$ is small.

\paragraph{Formal proof of Theorem~\ref{ch27:thm:echo}.} Under collaborative filtering with engagement maximization: the optimal $R^*$ maximizes $\sum_i \sum_j R^*_{ij} \cdot p_{ij}$, where $p_{ij}$ is the probability that user $i$ engages with content from source $j$. If $p_{ij}$ is decreasing in the ``distance'' between $i$ and $j$ (different group, different beliefs $\to$ less engagement), then $R^*_{ij}$ assigns more weight to within-group content. The induced $W^{I'}$ has $\lambda_k^{I'} \leq \lambda_k^I$.


\subsection*{Appendix 27B: The DeGroot Model with Algorithmic Intermediation}
\addcontentsline{toc}{subsection}{Appendix 27B: DeGroot with Algorithms}
\label{ch27:app:degroot_algo}

\paragraph{Modified DeGroot.} In Chapter~8, the DeGroot model had beliefs converging through a fixed $W^I$. With algorithmic intermediation, $W^I$ itself depends on beliefs:
\begin{equation}
\mathbf{b}_{t+1} = W^I(\mathbf{b}_t) \cdot \mathbf{b}_t,
\end{equation}
where $W^I(\mathbf{b}_t)$ is the algorithmically modified information matrix (the algorithm adjusts the feed based on current beliefs/engagement patterns). This is a \textit{nonlinear} dynamical system: the matrix changes with the state.

\paragraph{Fixed points.} A fixed point $\mathbf{b}^*$ satisfies $\mathbf{b}^* = W^I(\mathbf{b}^*)\mathbf{b}^*$: beliefs are a consensus of the information network that the beliefs themselves produced. This is a deeper form of the self-confirming trap: the algorithm creates the $W^I$ that produces the beliefs that the algorithm was designed to serve.

\paragraph{Multiple equilibria.} The nonlinear system can have multiple fixed points: a ``diverse'' equilibrium (where $W^I(\mathbf{b}^*)$ has high $N_{\text{eff}}$ because beliefs are varied, so the algorithm shows diverse content) and a ``polarized'' equilibrium (where $W^I(\mathbf{b}^*)$ has low $N_{\text{eff}}$ because beliefs are clustered, so the algorithm shows confirmatory content). The polarized equilibrium is absorbing under engagement-maximization: once beliefs cluster, the algorithm reinforces the clustering, preventing return to the diverse equilibrium.


%% =====================================================================
\subsection*{Structuralist Case Study: The Algorithmic $W^I$ in the Global South}
\addcontentsline{toc}{subsection}{Structuralist Case Study: Algorithmic $W^I$}
%% =====================================================================

This chapter's case study departs from the historical structuralist tradition to address a contemporary structural phenomenon. The algorithmic $W^I$ is arguably the most significant structural development since the colonial imposition of the racial classification system (Quijano, Chapter~18): both are externally imposed structures that shape $W^I$ for entire populations without their consent.

Three observations ground the structural analysis.

\textbf{First: the speed of $W^I$ transformation is unprecedented.} Bolivia went from zero social media users to 7 million in 12 years (2010--2022). In Chapter~23's terms: $W^I$ underwent a structural break --- not gradual evolution but a discrete regime change. The algorithmic $W^I$ did not grow out of the existing information structure; it replaced it. For urban Bolivians under 35, the algorithmic feed IS $W^I$ ($\omega_d = 0.50$, Chapter~21). This happened faster than any previous structural change in the book: faster than ISI, faster than nationalization, faster than the colonial imposition.

\textbf{Second: the structural effects are unmeasured.} Despite the speed and scale of the $W^I$ transformation, no developing country systematically monitors $N_{\text{eff}}$, $\lambda_k$, or $\Psi$ in the algorithmic $W^I$. The measurement infrastructure of Chapter~21 exists in principle but has not been deployed. We are flying blind: the most powerful $W^I$ operator in history is operating without structural monitoring.

\textbf{Third: the regulatory capacity is asymmetric.} The EU regulates algorithmic platforms (the Digital Services Act, 2022). The US debates regulation. The Global South --- where the structural effects may be largest (because $N_{\text{eff}}$ was already low and attribute constraints were already severe) --- has essentially no regulatory capacity. Bolivia's telecommunications regulator (ATT) has no mandate, tools, or data to assess the structural effects of Facebook's algorithm on Bolivia's $W^I$.

The structuralist response is clear: treat algorithmic $W^I$ as a structural variable that requires measurement (Chapter~21), monitoring (Chapter~23), and regulation (this chapter's Proposition~\ref{ch27:prop:regulation}). The alternative is to cede the formation of Bolivian preferences to an algorithm designed in Menlo Park --- the information-age equivalent of ceding the production structure to a colonial power.

\begin{center}
\begin{tabularx}{\textwidth}{LLL}
\toprule
\textbf{Structural parallel} & \textbf{Historical} & \textbf{Algorithmic} \\
\midrule
External control of $W$ & Colonial administration & Platform algorithms \\
Center-periphery & Trade asymmetry ($W^P$) & Information asymmetry ($W^I$) \\
Hegemonic formation & State media, schools & Engagement-maximizing feeds \\
Counter-hegemony & Community radio, movements & Public-interest algorithms? \\
Regulatory response & ISI, nationalization & Diversity mandates, $N_{\text{eff}}$ monitoring \\
\bottomrule
\end{tabularx}
\end{center}


%% =====================================================================
\subsection*{Further Reading}
\addcontentsline{toc}{subsection}{Further Reading}
%% =====================================================================

\textbf{Algorithmic curation and society:} Pariser (2011), \textit{The Filter Bubble}. Sunstein (2017), \textit{\#Republic: Divided Democracy in the Age of Social Media}. Zuboff (2019), \textit{The Age of Surveillance Capitalism}.

\textbf{Echo chambers and polarization:} Levy (2021), ``Social Media, News Consumption, and Polarization,'' \textit{AER}. Allcott et al.\ (2020), ``The Welfare Effects of Social Media,'' \textit{AER}. Bail et al.\ (2018), ``Exposure to Opposing Views on Social Media Can Increase Political Polarization,'' \textit{PNAS}.

\textbf{AI and development:} Acemoglu and Johnson (2023), \textit{Power and Progress: Our Thousand-Year Struggle over Technology and Prosperity}. World Bank (2019), \textit{World Development Report: The Changing Nature of Work}. For the Global South specifically: Arora (2019), \textit{The Next Billion Users: Digital Life Beyond the West}.

\textbf{Platform regulation:} EU Digital Services Act (2022). Stigler Committee on Digital Platforms (2019), \textit{Final Report}. For developing countries: Gurumurthy and Chami (2020), ``Platform Planet: Development in the Intelligence Economy,'' IT for Change.

\end{document}
